



% \vspace{-5mm}
% \subsubsection{Agent Specialization}
% \label{sec:specialization}

% 




\begin{table}[t]
\centering
\scriptsize
\setlength{\tabcolsep}{2pt}
\renewcommand{\arraystretch}{1.2}
% \vspace{-2mm}
\caption{Performance of agents under different conditions (using GPT-4o-mini). Precision (P), Recall (R), and F1 are shown.}

\label{tab:specialaiztion_result}
\vspace{-3mm}
\renewcommand{\arraystretch}{0.9}
\begin{tabular}{l|c|ccc|ccc|ccc|ccc}
\toprule
\textbf{Category} & \textbf{Regres.} & \multicolumn{3}{c|}{\textbf{Binary Class}} & \multicolumn{9}{c}{\textbf{Multiclass}} \\
\midrule
\multirow{2}{*}{\textbf{Agent Role}} 
& \multicolumn{1}{c|}{\textbf{Utility}} 
& \multicolumn{3}{c|}{\textbf{WiFi}} 
& \multicolumn{3}{c|}{\textbf{EU-IT}} 
& \multicolumn{3}{c|}{\textbf{Yelp}} 
& \multicolumn{3}{c}{\textbf{Volkert}} \\ 
\cmidrule(lr){2-2} \cmidrule(lr){3-5} \cmidrule(lr){6-8} \cmidrule(lr){9-11} \cmidrule(lr){12-14}
& \textbf{MAE} 
& \textbf{P} & \textbf{R} & \textbf{F1} 
& \textbf{P} & \textbf{R} & \textbf{F1} 
& \textbf{P} & \textbf{R} & \textbf{F1} 
& \textbf{P} & \textbf{R} & \textbf{F1} \\
\bottomrule
\multicolumn{14}{c}{\cellcolor{gray!20}\textbf{Role-based prompting}} \\
No Role        & 0.067 & 40.2 & 45.0 & 37.3  
& 32.0 & 38.0 & 35.0 
& 41.9 & 44.1 & 41.9 
& 67.3 & 66.5 & 65.3 \\

Data Scientist & 0.067 & 40.2 & 45.0 & 37.3 
& 32.3 & 37.9 & 34.8 
& 41.9 & 44.1 & 41.9 
& 67.3 & 66.5 & 65.3 \\

Researcher     & 0.067 & 40.2 & 45.0 & 37.3 
& 32.3 & 37.9 & 34.8 
& 42.1 & 44.0 & 41.9 
& 67.3 & 67.0 & 65.8 \\

Data Analyst   & 0.067 & 40.2 & 45.0 & 37.3 
& 33.3 & 37.9 & 34.8 
& 42.1 & 44.0 & 41.9
& 67.0 & 67.0 & 65.0 \\

Engineer       & 0.067 & 40.2 & 45.0 & 37.3 
& 32.0 & 38.0 & 35.0 
& 41.9 & 44.1 & 41.9 
& 67.3 & 66.5 & 65.3 \\

\multicolumn{14}{c}{\cellcolor{gray!20}\textbf{Planning-based conditioning}} \\
No Role        & 0.067 & 40.0 & 45.0 & 37.0 
& 32.0 & 38.0 & 35.0 
& 42.0 & 44.1 & 42.0
& 67.0 & 67.0 & 65.0 \\

Data Scientist & 0.067 & 40.0 & 45.0 & 37.0 
& 32.0 & 38.0 & 35.0 
& 42.0 & 44.1 & 42.0 
& 67.0 & 67.0 & 66.0 \\

Researcher     & 0.067 & 40.0 & 45.0 & 37.0
& 32.0 & 38.0 & 35.0 
& 42.0 & 44.0 & 42.0 
& 67.0 & 67.0 & 65.0 \\

Data Analyst   & 0.067 & 40.0 & 45.0 & 37.0
& 32.0 & 38.0 & 35.0 
& 42.0 & 44.0 & 42.0 
& 67.0 & 67.0 & 65.0 \\

Engineer       & 0.067 & 40.0 & 45.0 & 37.0
& 32.0 & 38.0 & 35.0 
& 42.0 & 44.0 & 42.0 
& 67.0 & 67.0 & 65.0 \\

\multicolumn{14}{c}{\cellcolor{gray!20}\textbf{Expert-guided conditioning}} \\
No Role        & 0.067 & 53.0 & 90.0 & 67.0 
& 33.6 & 41.4 & 36.8 
& 96.0 & 96.0 & 96.0 
& 94.8 & 94.7 & 94.6 \\

Data Scientist & 0.068 & 57.1 & 80.0 & 66.7 
& 37.3 & 44.8 & 40.3 
& 100.0 & 100.0 & 100.0 
& 95.0 & 95.0 & 95.0 \\

Researcher     & 0.068 & 53.0 & 90.0 & 67.0  
& 37.3 & 41.4 & 38.8 
& 100.0 & 100.0 & 100.0 
& 95.0 & 94.9 & 94.9 \\

Data Analyst   & 0.067 & 52.9 & 90.0 & 66.7 
& 37.3 & 41.4 & 38.8 
& 100.0 & 100.0 & 100.0 
& 94.7 & 94.6 & 94.5 \\

Engineer       & 0.068 & 50.0 & 90.0 & 64.3 
& 37.3 & 41.4 & 38.8 
& 100.0 & 100.0 & 100.0 
& 94.8 & 94.7 & 94.6 \\
\bottomrule
\end{tabular}
% \vspace{-5mm}
\end{table}

% Table~\ref{tab:specialaiztion_result} compares three textual conditioning strategies, role-based prompting, planning-based conditioning, and expert-guided conditioning, under a fixed model and execution pipeline, isolating the effect of conditioning on agent behavior. Role-based prompting alone produces negligible differences across roles: precision, recall, and F1 scores remain nearly identical across tasks, indicating that assigning professional identities such as \emph{Data Scientist}, \emph{Researcher}, or \emph{Engineer} does not reliably induce domain-specific reasoning. Planning-based conditioning further stabilizes behavior across roles, with metrics converging to similar values, but does not provide meaningful improvements.


% In contrast, expert-guided conditioning yields large and consistent gains across all tasks. Precision, recall, and F1 scores increase substantially for both binary and multiclass settings, with especially pronounced improvements on \textit{WiFi}, \textit{Yelp}, and \textit{Volkert}. These gains are consistent across roles, demonstrating that specialization is driven by explicit procedural guidance rather than identity framing. Overall, the results show that latent domain knowledge is present in the model but is only reliably activated through methodological instructions,\footnote{\vspace{-3mm} Example prompts are provided in the \uline{\textcolor{blue}{\href{https://github.com/CoDS-GCS/MAFBench/blob/main/supplementary_materials/APPENDICIES.pdf}{supplementary materials}}}.} 
% % (example prompts in Appendix~\ref{sec:Specialization_appendix})
% establishing specialization as a conditioning problem rather than a role assignment problem.



% \vspace{-1mm}
\subsubsection{Agent Specialization}
\label{sec:specialization}






\begin{table}[t]
\centering
\scriptsize
\setlength{\tabcolsep}{2pt}
\renewcommand{\arraystretch}{1.2}
% \vspace{-2mm}
\caption{Performance of agents under different conditions (using GPT-4o-mini). Precision (P), Recall (R), and F1 are shown.}

\label{tab:specialaiztion_result}
\vspace{-3mm}
\renewcommand{\arraystretch}{0.9}
\begin{tabular}{l|c|ccc|ccc|ccc|ccc}
\toprule
\textbf{Category} & \textbf{Regres.} & \multicolumn{3}{c|}{\textbf{Binary Class}} & \multicolumn{9}{c}{\textbf{Multiclass}} \\
\midrule
\multirow{2}{*}{\textbf{Agent Role}} 
& \multicolumn{1}{c|}{\textbf{Utility}} 
& \multicolumn{3}{c|}{\textbf{WiFi}} 
& \multicolumn{3}{c|}{\textbf{EU-IT}} 
& \multicolumn{3}{c|}{\textbf{Yelp}} 
& \multicolumn{3}{c}{\textbf{Volkert}} \\ 
\cmidrule(lr){2-2} \cmidrule(lr){3-5} \cmidrule(lr){6-8} \cmidrule(lr){9-11} \cmidrule(lr){12-14}
& \textbf{MAE} 
& \textbf{P} & \textbf{R} & \textbf{F1} 
& \textbf{P} & \textbf{R} & \textbf{F1} 
& \textbf{P} & \textbf{R} & \textbf{F1} 
& \textbf{P} & \textbf{R} & \textbf{F1} \\
\bottomrule
\multicolumn{14}{c}{\cellcolor{gray!20}\textbf{Role-based prompting}} \\
No Role        & 0.067 & 40.2 & 45.0 & 37.3  
& 32.0 & 38.0 & 35.0 
& 41.9 & 44.1 & 41.9 
& 67.3 & 66.5 & 65.3 \\

Data Scientist & 0.067 & 40.2 & 45.0 & 37.3 
& 32.3 & 37.9 & 34.8 
& 41.9 & 44.1 & 41.9 
& 67.3 & 66.5 & 65.3 \\

Researcher     & 0.067 & 40.2 & 45.0 & 37.3 
& 32.3 & 37.9 & 34.8 
& 42.1 & 44.0 & 41.9 
& 67.3 & 67.0 & 65.8 \\

Data Analyst   & 0.067 & 40.2 & 45.0 & 37.3 
& 33.3 & 37.9 & 34.8 
& 42.1 & 44.0 & 41.9
& 67.0 & 67.0 & 65.0 \\

Engineer       & 0.067 & 40.2 & 45.0 & 37.3 
& 32.0 & 38.0 & 35.0 
& 41.9 & 44.1 & 41.9 
& 67.3 & 66.5 & 65.3 \\

\multicolumn{14}{c}{\cellcolor{gray!20}\textbf{Planning-based conditioning}} \\
No Role        & 0.067 & 40.0 & 45.0 & 37.0 
& 32.0 & 38.0 & 35.0 
& 42.0 & 44.1 & 42.0
& 67.0 & 67.0 & 65.0 \\

Data Scientist & 0.067 & 40.0 & 45.0 & 37.0 
& 32.0 & 38.0 & 35.0 
& 42.0 & 44.1 & 42.0 
& 67.0 & 67.0 & 66.0 \\

Researcher     & 0.067 & 40.0 & 45.0 & 37.0
& 32.0 & 38.0 & 35.0 
& 42.0 & 44.0 & 42.0 
& 67.0 & 67.0 & 65.0 \\

Data Analyst   & 0.067 & 40.0 & 45.0 & 37.0
& 32.0 & 38.0 & 35.0 
& 42.0 & 44.0 & 42.0 
& 67.0 & 67.0 & 65.0 \\

Engineer       & 0.067 & 40.0 & 45.0 & 37.0
& 32.0 & 38.0 & 35.0 
& 42.0 & 44.0 & 42.0 
& 67.0 & 67.0 & 65.0 \\

\multicolumn{14}{c}{\cellcolor{gray!20}\textbf{Expert-guided conditioning}} \\
No Role        & 0.067 & 53.0 & 90.0 & 67.0 
& 33.6 & 41.4 & 36.8 
& 96.0 & 96.0 & 96.0 
& 94.8 & 94.7 & 94.6 \\

Data Scientist & 0.068 & 57.1 & 80.0 & 66.7 
& 37.3 & 44.8 & 40.3 
& 100.0 & 100.0 & 100.0 
& 95.0 & 95.0 & 95.0 \\

Researcher     & 0.068 & 53.0 & 90.0 & 67.0  
& 37.3 & 41.4 & 38.8 
& 100.0 & 100.0 & 100.0 
& 95.0 & 94.9 & 94.9 \\

Data Analyst   & 0.067 & 52.9 & 90.0 & 66.7 
& 37.3 & 41.4 & 38.8 
& 100.0 & 100.0 & 100.0 
& 94.7 & 94.6 & 94.5 \\

Engineer       & 0.068 & 50.0 & 90.0 & 64.3 
& 37.3 & 41.4 & 38.8 
& 100.0 & 100.0 & 100.0 
& 94.8 & 94.7 & 94.6 \\
\bottomrule
\end{tabular}
% \vspace{-5mm}
\end{table}

We evaluate how specialization mechanisms in multi-agent frameworks shape domain-specific reasoning behavior. Specifically, we isolate whether architectural conditioning through role assignment, planning interfaces, or procedural guidance governs performance. Using the unified specialization pipeline of MAFBench (Section~\ref{subsec:bench_specialization}), we fix the LLM backend, datasets from CatDB~\cite{catdb}, prompts, task structure, and evaluation metrics. We vary only the conditioning strategy applied to the agent. The three architectural variants correspond to role-based prompting, planning-based conditioning, and expert-guided procedural instructions.

Table~\ref{tab:specialaiztion_result} reports precision, recall, and F1 across regression, binary, and multiclass tasks. Role-based prompting produces nearly identical performance across all assigned professional identities, with no consistent improvement over the \emph{No Role} baseline. Planning-based conditioning further stabilizes behavior across roles but does not yield accuracy gains, indicating that introducing intermediate reasoning steps alone does not activate domain-relevant expertise. In contrast, expert-guided conditioning consistently and substantially improves performance on the classification datasets.Example prompts are provided in Appendix~\ref{sec:Specialization_appendix}.

These results show that specialization in LLM-based agents is governed by how frameworks inject task-specific reasoning structure into the execution pipeline, rather than by role identity alone. Role labels and generic planning interfaces fail to activate domain knowledge encoded in the LLM’s parameters. In contrast, explicit procedural instructions impose structured solution workflows that reliably improve performance. When designing multi-agent systems, specialization should be implemented through reasoning procedures embedded in the framework. Role naming or lightweight prompt modifications are insufficient, since execution structure directly determines robustness and effectiveness under fixed models.